\documentclass[a4paper, 12pt]{article}

\usepackage[utf8]{inputenc}
\usepackage[T1]{fontenc}
\usepackage[italian]{babel}
\addto\captionsitalian{\renewcommand{\abstractname}{Abstract}}

\usepackage[
  top=2.5cm, bottom=2.5cm,
  left=2.5cm, right=2.5cm
]{geometry}
\usepackage{fancyhdr}
\pagestyle{fancy}
    \fancyhf{}
    \fancyhead[L]{\small Laboratorio di Fisica Nucleare e Subnucleare 1}
    \fancyhead[R]{\small A.A. 2025/2026}
    \fancyfoot[C]{\thepage}
\fancypagestyle{plain}
{
  \fancyhf{}
  \fancyhead[L]{\small Laboratorio di Fisica}
  \fancyhead[R]{\small A.A. 2025/2026}
  \fancyfoot[C]{Università degli Studi di Torino\\A.A. 2025/2026}
}

\usepackage{soul}
\usepackage{xcolor}
\newcommand{\WIP}[1]{{\sethlcolor{cyan}\hl{#1}}}
\newcommand{\details}{{\sethlcolor{orange}\hl{inserire dettagli}}}


\usepackage{amsmath, amssymb}
\usepackage{physics}
\usepackage{siunitx}
\usepackage{xfrac}

\usepackage{graphicx}
\usepackage{adjustbox}
\usepackage{booktabs}
\usepackage{float}
\usepackage[
  justification=centering,
  labelfont={small,bf},
  font={small}
]{caption}
\usepackage{subcaption}
\usepackage{tikz}

\usepackage[compat=1.0.0]{tikz-feynman}
\usepackage[version=4,arrows=pgf-filled,
textfontname=sffamily,
mathfontname=mathsf]{mhchem}

\usepackage{parskip}
\usepackage{microtype}
\usepackage{hyperref}

\usepackage{appendix}
\usepackage[backend=biber, style=numeric]{biblatex}
\addbibresource{bibliografia.bib}


\title{\textbf{MISURA DELLA VITA MEDIA DEL MUONE}}
\author{Sara Barberis (sara.barberis939@edu.unito.it)
\\ Sara Ibba (sara.ibba815@edu.unito.it)
\\ Micaela Scamarcia (micaela.scamarcia@edu.unito.it)
\\ Mattia Valerio (mattia.valerio@edu.unito.it)
\\ Lorenzo Visca (lorenzo.visca846@edu.unito.it)}
\date{}

\begin{document}
\maketitle

\begin{center}
\vspace{60pt}
\begin{tikzpicture}
\begin{feynman}
  % Definizione vertici con coordinate esplicite
  \vertex (i1) at (-3,  0)  {\(\mu^-\)};
  \vertex (a)  at (-1,  0);
  \vertex (b)  at ( 1,  -1);
  \vertex (f1) at (3,  1.5)  {\(\nu_\mu\)};
  \vertex (f2) at ( 3,  -0.5)  {\(e^-\)};
  \vertex (f3) at ( 3, -1.5)  {\(\bar{\nu}_e\)};

  \diagram*{
    (i1) -- [fermion]                     (a),
    (a)  -- [boson, edge label=\(W^-\)]   (b),
    (a)  -- [fermion]                     (f1),
    (b)  -- [fermion]                     (f2),
    (b)  -- [anti fermion]                (f3),
  };
\end{feynman}
\end{tikzpicture}

\vspace{60pt}
\begin{tikzpicture}
\begin{feynman}
  % Definizione vertici con coordinate esplicite
  \vertex (i1) at (-3,  0)  {\(\mu^+\)};
  \vertex (a)  at (-1,  0);
  \vertex (b)  at ( 1,  -1);
  \vertex (f1) at (3,  1.5)  {\(\bar{\nu}_\mu\)};
  \vertex (f2) at ( 3,  -0.5)  {\(e^+\)};
  \vertex (f3) at ( 3, -1.5)  {\(\nu_e\)};

  \diagram*{
    (i1) -- [anti fermion]                (a),
    (a)  -- [boson, edge label=\(W^+\)]   (b),
    (a)  -- [anti fermion]                (f1),
    (b)  -- [anti fermion]                (f2),
    (b)  -- [fermion]                     (f3),
  };
\end{feynman}
\end{tikzpicture}
\end{center}


\newpage

\begin{abstract}
    L'esperienza ha avuto come obiettivo la misurazione della vita media dei muoni provenienti dai raggi cosmici, a tale scopo è stato realizzato un apparato sperimentale basato su rivelatori a scintillazione accoppiati a tubi fotomoltiplicatori e moduli di elettronica avanzata per la registrazione e l'analisi degli eventi.
    Il risultato da noi ottenuto è \WIP{[nostra stima della vita media del muone, dire velocemente se abbiamo accordo con quanto atteso]}.
    L'analisi temporale è stata affiancata da uno studio dello spettro energetico dei segnali uscenti dai rivelatori, evidenziando le correlazioni presenti tra il tipo di segnale osservato e l'energia depositata. \WIP{[rivedere in base a cosa troviamo effettivamente]}
  
\end{abstract}

\tableofcontents
\newpage

\section{Fisica dell'esperienza}
    L'interazione dei raggi cosmici primari con l'atmosfera terrestre porta alla produzione di sciami di adroni leggeri (principalmente $\pi^\pm$ e $\pi^0$, distribuiti in modo quasi uniforme con una leggera prevalenza di $\pi^+$ rispetto a $\pi^-$).
    \\
    Attraversando l'atmosfera una parte sempre maggiore dell'energia viene dispersa nel canale elettromagnetico tramite il processo di decadimento $\pi^0\to\gamma\gamma$, mentre i pioni carichi riescono a sostenere lo sciame per alcune generazioni prima di perdere energia e decadere.
    
    I pioni carichi decadono per interazione debole in un muone e un neutrino/antineutrino per conservazione del numero leptonico; i muoni così prodotti sono molto penetranti perché non partecipano ai processi di interazione forte e hanno una massa sufficiente a non rallentare in modo eccessivo per processi di bremsstrahlung, la loro vita media $\tau_\mu=2.197\unit{\micro\second}$ è sufficiente a farli arrivare fino al livello del mare, dove si registra un flusso di muoni di circa $130\unit{\hertz\per\meter\squared}$ rispetto ad un flusso totale di raggi cosmici di circa $180\unit{\hertz\per\meter\squared}$.
    
    I muoni decadono per interazione debole secondo i processi descritti dai seguenti diagrammi:
    
    \vspace{20pt}
    \begin{center}
    \begin{adjustbox}{width=0.6\textwidth}
    \begin{tikzpicture}
    \begin{feynman}
      % Definizione vertici con coordinate esplicite
      \vertex (i1) at (-3,  0)  {\(\mu^-\)};
      \vertex (a)  at (-1,  0);
      \vertex (b)  at ( 1,  -1);
      \vertex (f1) at (3,  1.5)  {\(\nu_\mu\)};
      \vertex (f2) at ( 3,  -0.5)  {\(e^-\)};
      \vertex (f3) at ( 3, -1.5)  {\(\bar{\nu}_e\)};
    
      \diagram*{
        (i1) -- [fermion]                     (a),
        (a)  -- [boson, edge label=\(W^-\)]   (b),
        (a)  -- [fermion]                     (f1),
        (b)  -- [fermion]                     (f2),
        (b)  -- [anti fermion]                (f3),
      };
    \end{feynman}
    \end{tikzpicture}
    \begin{tikzpicture}
    \begin{feynman}
      % Definizione vertici con coordinate esplicite
      \vertex (i1) at (-3,  0)  {\(\mu^+\)};
      \vertex (a)  at (-1,  0);
      \vertex (b)  at ( 1,  -1);
      \vertex (f1) at (3,  1.5)  {\(\bar{\nu}_\mu\)};
      \vertex (f2) at ( 3,  -0.5)  {\(e^+\)};
      \vertex (f3) at ( 3, -1.5)  {\(\nu_e\)};
    
      \diagram*{
        (i1) -- [anti fermion]                (a),
        (a)  -- [boson, edge label=\(W^+\)]   (b),
        (a)  -- [anti fermion]                (f1),
        (b)  -- [anti fermion]                (f2),
        (b)  -- [fermion]                     (f3),
      };
    \end{feynman}
    \end{tikzpicture}
    \end{adjustbox}
    \end{center}
    
    \vspace{20pt}
    Sperimentalmente vengono osservati gli eventi relativi ai muoni "lenti" $\left(T_\mu\sim100\unit{\mega\electronvolt}\right)$ che decadono all'interno del materiale attivo, che nel nostro caso è uno scintillatore composto da vetro scintillante, di seguito indicato con la sigla SG: i muoni che perdono tutta la loro energia all'interno del vetro si fermano e decadono, portando alla produzione di un elettrone (o positrone) che perde tutta la sua energia generando un secondo segnale all'interno dello scintillatore. Per conservazione dell'impulso nel processo di decadimento a tre corpi il leptone carico prodotto può assumere uno spettro continuo di energia con un endpoint $T_e=0.5m_\mu$, ottenuto quando il leptone carico viene emesso \textit{back-to-back} rispetto alla coppia di neutrini (che escono dal rivelatore senza produrre segnale).

    Nel caso dei muoni negativi, oltre al "puro" processo di decadimento è possibile anche osservare eventi di cattura del muone da parte dei nuclei atomici: se $\mu^-$ è abbastanza lento può essere catturato dal campo coulombiano del nucleo e scendere fino al livello fondamentale con l'emissione di raggi X ed elettroni Auger, una volta arrivato a tale livello si trova molto vicino al nucleo e può essere catturato per via radiativa ($\mu^-+p\to n+\nu_\mu+\gamma$) oppure non radiativa ($\mu^-+p\to n+\nu_\mu$), in questo caso si ha la formazione di un nucleo eccitato che decadendo può produrre un ulteriore segnale.
    \\
    La probabilità del processo e il tempo medio di cattura sono fortemente dipendenti dal materiale in cui si trova il muone, in materiali ad alto Z il tempo medio di cattura è nell'ordine di $10^2\unit{\nano\second}$.

\section{Apparato sperimentale}

    \subsection{Scintillatori e PMT}
    
        Il segnale relativo al processo studiato viene generato all'interno di uno scintillatore in vetro scintillante composto principalmente da \ce{SiO2} (42.5\%) e \ce{BaO3} (43.5\%), con centri di scintillazione di \ce{Ce2O3} (1.5\%).
        
        Lo scintillatore ha una densità $\rho=3.3\unit{\gram\per\centi\metre\cubed}$ e le seguenti dimensioni:

        \hl{nota: lunghezza sovrastimata perché c'è una specie di piastra alla fine, capire che incertezza mettere}
        
        \begin{table}[H]
            \centering
            \begin{tabular}{|c|c|c|}
            \toprule
            \multicolumn{3}{|c|}{\textbf{Scintillatore SG}} \\
            \midrule
            Lunghezza (\unit{\centi\metre}) & Larghezza (\unit{\centi\metre}) & Spessore (\unit{\centi\metre})\\
            \midrule
            $63.0\pm0.1$ & $17.0\pm0.1$ & $17.0\pm0.1$\\
            \bottomrule
            \end{tabular}
            \caption{Dimensioni di SG}
            \label{t:SG_dimensioni}
        \end{table}

        Lo spessore di massa dello scintillatore porta ad una perdita di energia media $\Delta E\approx110\unit{\mega\electronvolt}$ per particelle al minimo di ionizzazione che incidono verticalmente, quindi i muoni che incidono su SG con un'energia $\approx10^2\unit{\mega\electronvolt}$ perdono tutta la loro energia cinetica e decadono all'interno del materiale (nella pratica possono essere fermati anche muoni leggermente più energetici, che attraversando il materiale entrano nel regime $\sfrac{1}{\beta^2}$ della formula di Bethe-Bloch, e muoni ancora più energetici che attraversano lo scintillatore in diagonale).
        \\
        Lo scintillatore SG è accoppiato ad un tubo fotomoltiplicatore \details.

        La selezione di eventi potenzialmente "buoni" di SG viene effettuata tramite un sistema di trigger basato su due scintillatori plastici ($\rho\approx1\unit{\gram\per\centi\metre\cubed}$):

        \begin{itemize}
            \item uno scintillatore, indicato con S1 e posizionato sopra SG, il cui segnale viene messo in coincidenza con quello di SG per rimuovere gli eventi di buio di quest'ultimo;
            \item uno scintillatore, indicato con S2 e posizionato sotto SG, il cui segnale viene messo in anticoincidenza con quello di S1 ed SG per rimuovere gli eventi in cui il muone riesce a oltrepassare SG senza decadere.
        \end{itemize}
    
        La disposizione dei rivelatori è scelta basandosi sul fatto che il flusso di raggi cosmici non ha una distribuzione angolare isotropa, ma dipende dall'angolo $\theta$ rispetto all'asse zenitale:
    
        \begin{equation*}
            \dv{N}{\Omega}\propto\cos^2\theta
        \end{equation*}

        Durante l'esperienza è stato utilizzato anche un quarto scintillatore plastico più piccolo, indicato con S3, posizionato sopra S1 per migliorare la misura dell'efficienza di SG.

          \begin{table}[H]
            \centering
            \begin{tabular}{|c|c|c|c|}
            \toprule
            \midrule
            Scintillatore &Lunghezza (\unit{\centi\metre}) & Larghezza (\unit{\centi\metre}) & Spessore (\unit{\centi\metre})\\
            \midrule
            \textbf{S1} & $40.0\pm0.1$ & $11.0\pm0.1$ & $1.6\pm0.1$\\
            \textbf{S2} & $70.0\pm0.1$ & $16.0\pm0.1$ & $1.0\pm0.1$\\
            \textbf{S3} & $11.0\pm0.1$ & $10.0\pm0.1$ & $1.0\pm0.1$\\
            \bottomrule
            \end{tabular}
            \caption{Dimensioni degli scintillatori plastici}
            \label{t:S1_S2_S3_dimensioni}
        \end{table}     
        
        Lo spessore di massa degli scintillatori plastici porta ad una perdita di energia media inferiore ai $4\unit{\mega\electronvolt}$ per MIP verticali, questo valore è quindi trascurabile e non comporta un assorbimento significativo prima di SG.

        Tutti gli scintillatori plastici sono accoppiati a dei tubi fotomoltiplicatori \details XP2020 tramite delle guide di luce \textit{fishtail}.

    
    \subsection{Elettronica di acquisizione e analisi}
        \hl{non so se mettere il disegno completo dell'elettronica o l'elenco dei moduli e poi rimandare i singoli disegni alle varie fasi dell'esperienza}

        \hl{Se manteniamo l'elenco sarebbe carino mettere per ogni modulo il riferimento ai circuiti in cui viene utilizzato, ma è un po' sbattopoli}

        Di seguito è riportato l'elenco dei moduli di elettronica e degli strumenti utilizzati durante lo svolgimento dell'esperienza:
    
        \begin{itemize}
            \item \WIP{WIP}
            \item Scaler NIM \details
            \item Modulo di coincidenza \details
            \item Alimentatore ad alta tensione \details 
            \item Discriminatore \details
            \item Linear fan-out \details
            \item Impulsatore \details
            \item Oscilloscopio \details
            \item Multimetro \details
            \item Cavi segnale NIM \hl{standard connettore, ritardi}
            \item Impedenze di terminazione ("tappi") da $50\unit{\ohm}$ 
        \end{itemize}



     
    \newpage
\section{Inizializzazione e calibrazione dei rivelatori}

    \subsection{Tensione di lavoro degli scintillatori plastici}
        \hl{tensione di lavoro, curve di rate singolo, osservazione del segnale all'oscilloscopio}

    \subsection{Discriminazione dei segnali}
        \hl{Scelta soglie di discriminazione, scelta delle larghezze del segnale in uscita}

    \subsection{Temporizzazione dei segnali}
        \hl{Temporizzazione per curva di efficienza di SG, osservazione del time walk, S1 ed S2 risultano già temporizzati, spiegare perché S3 arriva dopo}

    \subsection{Tensione di lavoro di SG}
        \hl{tensione di lavoro, curve di rate singolo, osservazione del segnale all'oscilloscopio} \WIP{Curva di efficienza con le coincidenze 2-3, run notturna con S3 alla tensione scelta}

    \subsection{Curva di coincidenza}
        \WIP{Curva di coincidenza una volta scelta HV di SG, coincidenze accidentali, rate osservato vs atteso sia di segnale che di fondo confrontando i rate in singola visti lunedì}


\section{Implementazione del sistema di trigger}

\section{Misura della vita media del muone}

\section{Analisi ADC}

\section{Risultati ottenuti}

\section{Conclusioni}

\printbibliography

\newpage
\begin{appendices}
  \section{a}
\end{appendices}

\end{document}

%nota per il futuro: ricordiamoci di uniformare la notazione usata per descrivere il segnale es. altezza/ampiezza, lunghezza/durata
%nota da Visca: se scrivo da fisso ho la tastiera inglese che non ha gli accenti, ogni tanto scrivo parole in inglese a caso because non ho voglia di cercare gli accenti poi sistemo da portatile
%sto usando /WIP per le cose che dovremo aggiungere alla fine, /hl per le cose che non avevo voglia di fare adesso